% Part: first-order-logic
% Chapter: axiomatic-deduction
% Section: provability

% verification of properties of provability needed for maximally
% consistent sets in the completeness chapter.

\documentclass[../../../include/open-logic-section]{subfiles}

\begin{document}

\iftag{FOL}
      {\olfileid{fol}{axd}{prv}}
      {\olfileid{pl}{axd}{prv}}

\olsection{可推导性的性质}

\begin{prop}[单调性] 若 $\Gamma \subseteq \Delta$ 且 $\Gamma \Proves !A$，
则亦有 $\Delta \Proves !A$。\end{prop}

\begin{proof} 任何有限的 $\Gamma_0 \subseteq \Gamma$ 也是 $\Delta$ 的有限子集，
因此 $\Gamma_0 \Proves !A$ 的一个推导同时也表明 $\Delta \Proves !A$。\end{proof}

\begin{prop}\ollabel{prop:provability} $\Gamma \Proves !A$ 当且仅当
$\Gamma\cup \{\lnot!A \}$ 是不一致的。\end{prop}

\begin{prop} \begin{enumerate}
\item \ollabel{prop:provability-contr} 若 $\Gamma \Proves !A$ 且 $\Gamma
\cup \{ !A\} \Proves \lfalse$，则 $\Gamma$ 是不一致的。

\item \ollabel{prop:provability-lnot} 若 $\Gamma \cup \{!A\}
\Proves \lfalse$，则 $\Gamma \Proves \lnot !A$。

\item \ollabel{prop:provability-exhaustive} 若 $\Gamma \cup \{!A\}
\Proves \lfalse$ 且 $\Gamma \cup \{\lnot !A\} \Proves
\lfalse$，则 $\Gamma \Proves \lfalse$。

\tagitem{prvOr}{\ollabel{prop:provability-lor-left} 若 $\Gamma \cup \{!A\}
\Proves \lfalse$ 且 $\Gamma \cup \{!B\} \Proves
\lfalse$，则 $\Gamma \cup \{!A \lor !B\} \Proves \lfalse$。}{}

\tagitem{prvOr}{\ollabel{prop:provability-lor-right} 若 $\Gamma
\Proves !A$ 或 $\Gamma \Proves !B$，则 $\Gamma
\Proves !A \lor !B$。}{}

\tagitem{prvAnd}{\ollabel{prop:provability-land-left} 若 $\Gamma
\Proves !A \land !B$ 则 $\Gamma \Proves !A$ 且
$\Gamma \Proves !B$。}{}

\tagitem{prvAnd}{\ollabel{prop:provability-land-right} 若 $\Gamma
\Proves !A$ 且 $\Gamma \Proves !B$，则 $\Gamma
\Proves !A \land !B$。}{}

\tagitem{prvIf}{\ollabel{prop:provability-mp} 若 $\Gamma \Proves
!A$ 且 $\Gamma \Proves !A \lif !B$，则 $\Gamma
\Proves !B$。}{}

\tagitem{prvIf}{\ollabel{prop:provability-lif} 若 $\Gamma \Proves
\lnot !A$ 或 $\Gamma \Proves !B$，则 $\Gamma \Proves
!A \lif !B$。}{} \end{enumerate} \end{prop}

\begin{proof}
\begin{enumerate}
\item 假设 $\Gamma_0 \cup \{!A\} \Proves \lfalse$。

\begin{derivation}
1. & $\Gamma_0 \cup \{!A\} \Proves \lfalse$ & 假设 \\
2. & $\Gamma_1 \Proves !A$ & 假设\\
3. & $\Gamma_0 \cup \Gamma_1 \Proves \lfalse$ & Cut \\
\end{derivation}

由于 $\Gamma_0 \subseteq \Gamma$ 且 $\Gamma_1 \subseteq \Gamma$，
$\Gamma_0 \cup \Gamma_1 \subseteq \Gamma$，因此 $\Gamma \Proves \lfalse$。

\item 假设 $\Gamma_0 \cup\{\lnot !A\} \Proves \lfalse$。

\begin{derivation}
1. & $\Gamma_0 \cup\{\lnot !A\} \Proves \lfalse$ & 假设 \\
2. & $\Gamma_0 \Proves \lnot !A \lif \lfalse$ & 演绎定理\\
3. & $\Gamma_0 \Proves (\lnot !A \lif \lfalse) \lif !A$ & Ax0\\
4. & $\Gamma_0 \Proves !A$ & MP 2, 3 \\
\end{derivation}

\item 假设 $\Gamma \cup \{ !A \} \Proves \lfalse$ 且 $\Gamma_1 \cup \{\lnot !A \} \Proves \lfalse$。

\begin{derivation}
1. & $\Gamma_0 \cup\{!A\} \Proves \lfalse$ & 假设 \\
2. & $\Gamma_0 \Proves !A \lif \lfalse$ & 演绎定理 \\
3. & $\Gamma_1 \cup \{\lnot !A \} \Proves \lfalse$ & 假设\\
4. & $\Gamma_0 \Proves (!A \lif \lfalse) \lif \lnot !A$ & Ax0\\
5. & $\Gamma_0 \Proves \lnot !A$ & MP 2, 4\\
6. & $\Gamma_0 \cup \Gamma_1 \Proves \lfalse$ & Cut 3, 5 \\
\end{derivation}

由于 $\Gamma_0 \subseteq \Gamma$ 且 $\Gamma_1 \subseteq \Gamma$，
$\Gamma_0 \cup \Gamma_1 \subseteq \Gamma$。因此 $\Gamma \Proves \lfalse$。

% prop:provability-lor-left
\tagitem{defOr}{}{
\iftag{probOr}{练习。}{练习。}}

% prop:provability-lor-right
\tagitem{defOr}{}{ \iftag{probOr}{练习。}{
假设 $\Gamma_0 \Proves !A$。

\begin{derivation}
1. & $\Gamma_0 \Proves !A$ & 假设 \\
2. & $\Gamma_0 \Proves !A \lif (!A \lor !B)$ & Ax0\\
3. & $\Gamma_0 \Proves !A \lor !B$ & MP 1, 2 \\
\end{derivation}

因此 $\Gamma \Proves !A \lor !B$。当 $\Gamma \Proves
!B$ 时的证明类似。}}

% prop:provability-land-left
\tagitem{defAnd}{}{ \iftag{probAnd}{练习}{
假设 $\Gamma_0 \Proves !A \land !B$

\begin{derivation}
1. & $\Gamma_0 \Proves !A \land !B$ & 假设 \\
2. & $\Gamma_0 \Proves (!A\land !B) \lif !A $ & Ax0\\
3. & $\Gamma_0 \Proves !A \lor !B$ & MP 1, 2 \\
\end{derivation}

因此，$\Gamma \Proves !A$。一个类似的推导表明
$\Gamma \Proves !B$。}}

% prop:provability-land-right
\tagitem{defAnd}{}{\iftag{probAnd}{练习。}{练习。}}

% prop:provability-mp
\tagitem{defIf}{}{ \iftag{probIf}{练习。}{ 假设
$\Gamma_0 \Proves !A$ 且 $\Gamma_0 \Proves !A \lif !B $。

\begin{derivation}
1. & $\Gamma_0 \Proves !A$ & 假设 \\
2. & $\Gamma_0 \Proves !A \lif !B $ & 假设\\
3. & $\Gamma_0 \Proves !B$ & MP 1, 2\\
\end{derivation}

由于 $\Gamma_0 \cup \Gamma_1 \subseteq \Gamma$，这意味着 $\Gamma
\Proves !B$。}}

% prop:provability-lif
\tagitem{defIf}{}{\iftag{probIf}{练习。}{首先
假设 $\Gamma \Proves \lnot !A$。

\begin{derivation}
1. & $\Gamma_0 \Proves \lnot !A$ & 假设 \\
2. & $\Gamma_0 \Proves \lnot !A \lif (!A \lif !B) $ & Ax0\\
3. & $\Gamma_0 \Proves !A \lif !B$ & MP 1,
2\\ \end{derivation}

$\Gamma \Proves !B$ 的证明类似：

\begin{derivation}
1. & $\Gamma_0 \Proves !B$ & 假设 \\
2. & $\Gamma_0 \Proves !B \lif (!A \lif !B) $ & Ax0\\
3. & $\Gamma_0 \Proves !A \lif !B$ & MP 1, 2\\
\end{derivation}}}

\end{enumerate}
\end{proof}

\begin{probtag}{probOr,probAnd,probIf} 完成
\olref[fol][axd][prv]{prop:provability} 的证明。
\end{probtag}

\begin{thm}[泛化定理] \ollabel{thm:Generalization} 若 $\Gamma \Proves
!A$ 且 $x$ 在 $\Gamma$ 的任何公式中均不自由，则 $\Gamma
\Proves \lforall[x][!A]$。
\end{thm}

\begin{proof} 通过对 $\mathsf{Thm}_\Gamma$ 进行归纳来证明：若 $!A$ 是公理，
则 $\lforall[x][!A]$ 也是公理。若 $!A\in \Gamma$，则 $x$ 在 $!A$ 中不自由，故
$!A \lif \lforall[x][!A]$ 是一条公理（\textbf{Ax3}），并且 $\Gamma
\Proves \lforall[x][!A]$ 可由 MP 规则得到。现假设 $!A$ 由 MP 规则得出，即由
$\Gamma \Proves !B$ 且 $\Gamma \Proves !B \lif !A$ 而得。由归纳
假设，$\Gamma \Proves \lforall[x][!B]$ 且 $\Gamma \Proves
\lforall[x][( !B \lif !A)]$。但由 \textbf{Ax2} 有 \[ \Gamma \Proves
\lforall[x][( !B \lif !A)] \lif (\lforall[x][ !B] \lif \lforall[x][!A]), \]
两次应用 MP 规则即得所证的 $\Gamma \Proves \lforall[x][!A]$。 \end{proof}

\begin{thm}\ollabel{thm:WeakGen} （\emph{弱常项泛化定理}）
若 $\Gamma \Proves !A$ 且常项符号 $c$ 不在 $\Gamma$ 中出现，
则存在一个不在 $!A$ 中出现的变元 $x$，使得 $\Gamma
\Proves \lforall[x][\Subst{!A}{x}{c}]$，并且该证明不涉及 $c$。
\end{thm}

\begin{proof} 令 $!A_1,\ldots,!A_n$ 为从 $\Gamma$ 推出 $!A$ 的一个证明，故
$!A=!A_n$。选取一个不在 $!A_1,\ldots,!A_n$ 中出现的变元 $x$，并考虑
新的序列：$\Subst{!A_1}{x}{c},\ldots,\Subst{!A_n}{x}{c}.$ 该
序列是 $\Subst{!A}{x}{c}$ 从 $\Gamma$ 的一个证明。事实上，对每个
$i=1,\ldots,n$： \begin{itemize} \item 若 $!A$ 是公理，则 $\Subst{!A_i}{x}{c}$ 也是公理； \item 若 $!A_i \in \Gamma$，则由于 $c$ 不在
$\Gamma$ 中，我们有 $\Subst{!A_i}{x}{c} = !A_i \in \Gamma$； \item 若 $!A_i$
是由 $!A_j$ 和 $!A_j \lif !A_i$ 通过 MP 规则得到，则 $\Subst{!A_i}{x}{c}$
可由 MP 规则从 $\Subst{!A_j}{x}{c}$ 和 $\Subst{(!A_j \lif !A_i)}{x}{c}
= \Subst{!A_j}{x}{c} \lif \Subst{!A_i}{x}{c}$ 得到。\end{itemize} 显然，
常项符号 $c$ 不再出现在新序列中。现在令
$\Gamma'$ 由出现在 $\Subst{!A}{x}{c}$ 推导中的那些 $\Gamma$ 的公式组成；那么 $x$ 在 $\Gamma'$
的任何公式中均不自由，且 $\Gamma' \Proves \Subst{!A}{x}{c}$。由
泛化定理（\olref{thm:Generalization}）得 $\Gamma'
\Proves \lforall[x][\Subst{!A}{x}{c}]$，再由单调性得 $\Gamma \Proves
\lforall[x][\Subst{!A}{x}{c}]$，如所求。
\end{proof}

\begin{explain}
我们的目标是用更弱的要求——$x$ 在 $!A$ 中不自由且对 $c$ 在 $!A$ 中自由——来
替换 \olref{thm:WeakGen} 中 $x$ 完全不在 $!A$ 中出现这一要求。显然这可以通过
改变约束变元来实现——而这正是我们首先要证明的。
\end{explain}

\begin{lem}
\ollabel{lem:free-for}
若 $x$ 对 $c$ 在 $!A$ 中自由且 $y$ 在 $!A$ 中不自由，则 $x$
对 $y$ 在 $\Subst{!A}{y}{c}$ 中自由。
\end{lem}

\begin{proof} 若 $x$ 对 $y$ 在 $\Subst{!A}{y}{x}$ 中不自由，则
$\Subst{!A}{y}{c}$ 中 $y$ 的某个自由出现在量词 $\lforall[x]$ 的辖域内。
但由假设 $y$ 在 $!A$ 中不自由，所以所有这些出现都源自代入 $\Subst{}{y}{c}$。因此，
$c$ 的某个出现在 $\lforall[x]$ 的辖域内，从而 $x$ 对 $c$ 在 $!A$ 中不自由。
\end{proof}

\begin{lem}[改变约束变元]
\ollabel{lem:ChangeBdVar}
若 $x$ 和 $y$ 在 $!A$ 中均不自由，且两者都对 $c$ 在 $!A$ 中自由，
则 $\Proves \lforall[x][\Subst{!A}{x}{c}] \equiv
\lforall[y][\Subst{!A}{y}{c}]$（且证明不涉及 $c$）。
\end{lem}

\begin{proof}
由假设，$y$ 对 $c$ 在 $!A$ 中自由且 $x$ 在 $!A$ 中不自由，故由
前一引理 \olref{lem:free-for}，$y$ 对 $x$ 在 $\Subst{!A}{x}{c}$ 中自由。由此可知
$\lforall[x][\Subst{!A}{x}{c} \lif \Subst{!A}{x}{c} \Subst{}{y}{x}]$
是一条公理（\textbf{Ax1}）。由于 $x$ 在 $!A$ 中不自由，我们有
$\Subst{!A}{x}{c} \Subst{}{y}{x} = \Subst{!A}{y}{c}$，所以
\[
\Proves \lforall[x][\Subst{!A}{x}{c}] \lif \Subst{!A}{y}{c},
\]
并且由演绎定理得 $\lforall[x][\Subst{!A}{x}{c}] \Proves
\Subst{!A}{y}{c}$。由于 $y$ 在 $!A$ 中不自由，它在
$\lforall[x][\Subst{!A}{x}{c}]$ 中也不自由，故由泛化定理
$\lforall[x][\Subst{!A}{x}{c}] \Proves \lforall[y][\Subst{!A}{y}{c}]$，
再次应用演绎定理得 $\Proves
\lforall[x][\Subst{!A}{x}{c}] \lif \lforall[y][\Subst{!A}{y}{c}]$。
$\Proves \lforall[y][\Subst{!A}{y}{c}] \lif \lforall[x]
[\Subst{!A}{x}{c}]$ 的证明完全对称，因此结论由重言式可得。
\end{proof}

\begin{thm}[强常项泛化定理]
\ollabel{thm:StrongGen}
若 $\Gamma \Proves !A$，常项符号 $c$ 不在 $\Gamma$ 中出现，
且 $x$ 在 $!A$ 中不自由但对 $c$ 在 $!A$ 中自由，
则 $\Gamma \Proves \lforall[x][\Subst{!A}{x}{c}]$。
\end{thm}

\begin{proof}
由于 $\Gamma \Proves !A$ 且 $c$ 不在 $\Gamma$ 中，由
弱常项泛化定理知，存在一个不在 $!A$ 中出现的变元 $y$，使得
$\Gamma \Proves \lforall[y][\Subst{!A}{y}{c}]$。由于 $y$ 完全不在 $!A$ 中出现，
它在 $!A$ 中不自由，并且对 $c$ 在 $!A$ 中也自由；
又由假设 $x$ 在 $!A$ 中不自由且对 $c$ 在 $!A$ 中自由，
则改变约束变元的条件得到满足，故
$\Proves \lforall[x][\Subst{!A}{x}{c}] \equiv
\lforall[y][\Subst{!A}{y}{c}]$，进而得 $\Gamma \Proves
\lforall[x][\Subst{!A}{x}{c}]$。
\end{proof}

\begin{thm}
\ollabel{thm:provability-quantifiers}
\begin{tagenumerate}{prvEx,prvAll}
\tagitem{prvEx}{若 $\Gamma \Proves !A(t)$ 则 $\Gamma \Proves
  \lexists[x][!A(x)]$。}{}
\tagitem{prvAll}{若 $\Gamma \Proves \lforall[x][!A(x)]$ 则 $\Gamma
  \Proves !A(t)$。}{}
\end{tagenumerate}
\end{thm}

\begin{proof}
\begin{tagenumerate}{prvEx,prvAll}
%%Need axioms for exists
\tagitem{prvEx}{练习。}{}
\tagitem{prvAll}{假设 $\Gamma_0 \Proves \lforall[x][!A(x)]$。

\begin{derivation}
1. & $\Gamma_0 \Proves \lforall[x][!A(x)]$ & 假设 \\
2. & $\Gamma_0 \Proves \lforall[x][!A(x)] \lif \Subst{!A}{t}{x}$ & Ax1 \\
3. & $\Gamma_0 \Proves \Subst{!A}{t}{x}$ & MP 1, 2 \\
\end{derivation}}{}

\end{tagenumerate}
\end{proof}

\end{document}